\scanpage{56}
\mspage{12}
% archive sheet number: 55
\deleted{\unsure{(A)}}\inserted{(B)} \uline{Opération de \deleted{\unsure{affinement et de}}\inserted{sur-}marquage}

\deleted{On peut le définir \inserted{\unsure{en termes de}} par la donnée, d'\unsure{un} \unsure{sur chaque|jeu de} \unsure{int.|points} \unsure{en nombre donné} de \unsure{contour} $R$, $\unsure{R\times g}$ $\widetilde\Sigma\setminus R$. Soit $\unsure{\widetilde\Sigma_*}$ $A^+_{\unsure{0}}$ \unsure{et} l'ens \unsure{des}}

Soit $A^+$ l'ens des arêtes \inserted{(libres ou liées)} \inserted{de $\unsure{\Gamma|H}$} telles que $r(a)>0$ (``effectivement marquées'') \deleted{et $B$ l'ens des}, correspondant à des $C\in\pi_0(\widetilde\Sigma)$ telles que $C\cap R\ne\varnothing$, soit $B^+$ l'ens des comp. connexes de
\[
\widetilde\Sigma^+\setminus R\cap\widetilde\Sigma^+
\]
(lesquelles sont des int. ouverts), un surmarquage est défini par la donnée d'une application
\[
B^+\xrightarrow{r'}\mathbb N,
\]
% slanted
\margin{\deleted{\unsure{arêtes}} cercles \deleted{\unsure{effectifs}}\inserted{\unsure{effectivement}} marqués, \inserted{\unsure{corr.}} aux \uline{arêtes} sur $\widetilde\Sigma$ \unsure{que} (\uline{NB} elle est \unsure{réunion} des cercles \unsure{eff.|effectivement} marqués, et \unsure{non} des intervalles ouverts)}
et peut se décrire \deleted{par}\inserted{comme} l'opération qui \deleted{\unsure{qui}} sur chaque \uline{arête}\inserted{(\unsure{connexe|conn.})} de $\widetilde\Sigma$, correspondant à un $b\in B^+$, $r'(b)$ nouveaux sommets de \uline{rigidification}. Cette fois-ci, \deleted{\unsure{le point}} il y a des choix arbitraires -- néanmoins l'opération donne un objet de $S\mathscr T\unsure{\mathbb D|D}^+$, défini à isom unique près -- \inserted{soit $(X,\widetilde\Sigma,R_1)$} muni bien entendu d'un ens marqué $R'\subset R_1$, tel que l'opération inverse de gommage de $R'$ redonne l'objet de départ $(X,\widetilde\Sigma,R)$.
